\documentclass[12pt]{article}
\usepackage[utf8]{inputenc}
\usepackage[T1]{fontenc}
\usepackage{hyperref}
\title{Iau Language Resources for K-12 Curriculum Development}
\author{Compiled by Claude Sonnet 5 (claude-sonnet-5) for Andr\'as Kornai}
\date{}
\begin{document}
\maketitle

\section{Introduction}
Iau (ISO 639-3 tmu; Glottocode iauu1242) is a Lakes Plain language of Papua, Indonesia, with roughly 1{,}500 to 2{,}100 speakers. Most speakers are reported as monolingual. Searches found very little material for Iau. Other Lakes Plain languages such as Fayu turned up in the same searches, and their resources are not listed here.

\section{Primary Educational Resources}
No primers, textbooks or school programmes for Iau were found in search results.

\section{Linguistic and Reference Materials}
\begin{itemize}
\item Wikipedia overview of the Iau language, covering its tone system and dialects (Foi, Turu, Edopi, Iau proper): \url{https://en.wikipedia.org/wiki/Iau_language}
\item ASJP word list for Iau: \url{https://asjp.clld.org/languages/IAU}
\end{itemize}

\section{Religious and Community Resources}
No Bible translation or community organisation for Iau was found in search results.

\section{Digital and Audio Resources}
No audio, video or software resources for Iau were found in search results.

\section{Curriculum Development Context}
Very little published material for Iau was found. Curriculum work would need new fieldwork and an orthography, with attention to the very complex tone system. Further sources to check are SIL, Ethnologue and the literature on Lakes Plain languages.
\end{document}